\documentclass[11pt,a4paper]{article}
\usepackage[T1]{fontenc}
\usepackage{lmodern}
\usepackage[margin=19mm]{geometry}
\usepackage{graphicx,booktabs,xcolor,amsmath,hyperref}
\hypersetup{colorlinks=true,urlcolor=blue}
\setlength{\parindent}{0pt}
\setlength{\parskip}{6pt}
\setlength{\tabcolsep}{8pt}
\renewcommand{\arraystretch}{1.13}
\newcommand{\titulo}[1]{\textbf{\textcolor{blue!45!black}{#1}}\par}
\begin{document}
\fontsize{11}{14}\selectfont

\titulo{UTFPR · CSI30 · Sistemas Inteligentes · 2026/2}
\textbf{Tarefa 3 — Agrupamento de vítimas no VictSim3}\\
Julia Kamilly de Oliveira · RA 2588005 · Prof. Cesar Augusto Tacla

\titulo{1. Experimento básico com posição espacial}

O mapa explorado tem 5.625 células e 410 vítimas com ID único. A grade é
$75\times75$, com base absoluta $(30,40)$, 555 paredes e terreno de dificuldade variável.
A análise usa as coordenadas relativas à base e somente as células com ID não negativo.
Para cada $K$, aplicou-se K-means euclidiano com inicialização k-means++, 50 reinícios e semente 42.

\textbf{Tabela 1.} Variação de $K$ com somente $x_{rel}$ e $y_{rel}$.
\begin{center}
\begin{tabular}{rrr}
\toprule
$K$ & SSE (células$^2$) & Silhueta média \\
\midrule
\input{../outputs/basico_linhas.tex}
\end{tabular}
\end{center}

\begin{center}
\includegraphics[width=.495\linewidth]{../outputs/graficos/sse.pdf}
\includegraphics[width=.495\linewidth]{../outputs/graficos/silhueta.pdf}\\
\textbf{Figura 1.} SSE por $K$.\hfill\textbf{Figura 2.} Silhueta por $K$.
\end{center}

O enunciado menciona uma raiz quadrada e nomeia a métrica da tabela como SSE. Foi usada
a inércia do K-means, como no \href{https://github.com/tacla/Sistemas-Inteligentes/blob/187c1f0ca691d40a880616ef0517bf8134585687/Clustering/K_MEANS_An\%C3\%A1lises.ipynb}{exemplo de Tacla}, $\mathrm{SSE}=\sum_i\lVert x_i-\mu_{c_i}\rVert_2^2$.
A média quadrática por ponto é $\mathrm{SSE}/N$ e o RMSE correspondente é
$\sqrt{\mathrm{SSE}/N}$. A silhueta global é a média das silhuetas individuais com distância
euclidiana.

O cotovelo está em $K=4$, maior distância perpendicular à corda que liga $K=2$ e $K=12$
com os dois eixos escalados em $[0,1]$ apenas para localizar o ponto. O ganho de SSE
cai de 51.011,24 ($3\to4$) para 17.326,75 ($4\to5$). A silhueta atinge o máximo
em $K=5$ (0,4815), próxima da observada em $K=4$ (0,4753).

\clearpage
\titulo{2. Estratégia de agrupamento para o socorro}

\textbf{Predições.}
O mapa de entrada contém \texttt{tri} e \texttt{sobr} iguais a $-1$ nas 5.625 células.
Nas 410 vítimas, aplicou-se a rede da T1 com seu scaler interno e a rede da T2
com seu scaler externo. A ordem foi idade, fc, fr, pas, spo2, temp, pr, sg, fx e queim.
As escolhas seguiram os testes cegos anteriores (F1 macro 0,93020 e RMSE 0,04913).
Foram previstos 100 verdes, 106 amarelos, 108 vermelhos e 96 pretos.
A sobrevivência variou de 0,1656 a 0,9582, sem ajustes aos limites.

\textbf{Representação.}
As coordenadas absolutas divididas por 74 preservam a escala comum dos eixos.
A sobrevivência prevista $p$ já está em $[0,1]$ e entra como $0{,}45p$.
Uma diferença de 0,1 em $p$ equivale a 3,33 células no cálculo da distância.
A posição favorece lotes próximos e a sobrevivência diferencia perfis dentro do território.
A triagem tem ordem de prioridade e a classe preta requer tratamento próprio. Seus códigos
carecem de intervalos clínicos iguais. As medianas de $p$ por classe
são 0,938, 0,784, 0,502 e 0,251, da verde à preta. A triagem orienta a fila futura.
A idade entra como restrição de elegibilidade na atribuição.

\begin{center}
\includegraphics[height=73mm]{../outputs/graficos/nuvem.pdf}\\
\textbf{Figura 3.} Nuvem no espaço efetivo do agrupamento.
\end{center}

\textbf{Geometria.}
A nuvem contém vazios e regiões adensadas, com sobrevivências variadas
nas mesmas regiões. As correlações de Spearman de $p$ com $x$ e $y$ são $-0{,}057$ e
$0{,}074$. A distância ao quinto vizinho (incluindo o próprio ponto) tem mediana 0,086
e máximo 0,220, sinal de pontos mais isolados. A projeção sugere uma distribuição irregular,
sem evidência de grupos naturais isolados com contornos não convexos. Os centroides
fornecem setores de trabalho aproximados e podem cruzar paredes.

\textbf{Algoritmo.}
O DBSCAN permite contornos não convexos. Com cinco amostras mínimas e
$\varepsilon$ nos quartis da distância do quinto vizinho, produziu 16, 14 e 10 grupos
e 236, 116 e 45 casos de ruído, respectivamente. As coberturas foram 42,4\%, 71,7\%
e 89,0\%. Adotou-se K-means por sua partição de cobertura total e por permitir controlar
o número de lotes de trabalho. Nenhuma vítima foi excluída.

\clearpage
\titulo{3. Parametrizações e coordenação proposta}

\textbf{Tabela 2.} K-means no espaço da estratégia. O raio RMS mede a dispersão espacial
em células relativa ao centroide espacial de cada grupo.
\begin{center}
\begin{tabular}{rrrrrr}
\toprule
PAR & $K$ & SSE & Silhueta & Vítimas/lote & Raio RMS \\
\midrule
\input{../outputs/estrategia_linhas.tex}
\end{tabular}
\end{center}

A configuração escolhida fixa $K=6$. Em relação a $K=4$, reduz o raio RMS em cerca de 20\%,
com perda de 0,0248 na silhueta. Oito lotes equilibram melhor as contagens, mas o ganho
adicional de 0,98 célula no raio vem com perda de 0,0637 na silhueta em relação a seis.
Com $K=6$, os pesos 0,25 e 0,70 tiveram índices Rand ajustados de 0,938 e 0,533 em relação
a 0,45. A redução de peso mantém uma partição semelhante. O aumento reorganiza
consideravelmente os setores. SSEs de espaços ponderados de formas diferentes têm escalas distintas.

\textbf{Tabela 3.} Composição dos seis lotes e localização em relação à base.
\begin{center}
\begin{tabular}{rrrrrr}
\toprule
Lote & Total & Idade $\leq12$ & Centro $x_{rel}$ & Centro $y_{rel}$ & Dist. média \\
\midrule
\input{../outputs/grupos_linhas.tex}
\end{tabular}
\end{center}

\textbf{Atribuição futura.}
Os agentes 1 e 2 atendem qualquer idade e terão capacidade reservada para as 106
vítimas de até 12 anos. Agentes 3, 4 e 5 recebem somente vítimas com idade $>12$.
Em lotes mistos, dividir a fila em listas elegíveis por vítima. Por exemplo, no lote 1,
30 das 115 vítimas exigem agentes 1 ou 2 e as 85 restantes podem ser distribuídas
entre os cinco. O sexto lote dá margem para colocar trabalho em fila. Priorizar trabalho
por START (vermelho, amarelo, verde), revisar casos pretos separadamente.
Dentro de cada classe, ordenar pelo custo de deslocamento e antecipar a menor
sobrevivência estimada quando os custos forem semelhantes.
As previsões são estimativas. O custo de rota considera posições, relevo e paredes.

\textbf{Coordenação futura.}
Usar um coordenador com fila de IDs únicos. Trabalhar em paralelo em setores diferentes
ou em subconjuntos disjuntos de um lote grande. Reservar cada ID para um agente de
cada vez até conclusão ou liberação e redistribuir pendências conforme carga restante,
tempo estimado de atendimento e custo desde a posição atual ou a base.
A distância média da tabela é euclidiana e serve de referência inicial.
A execução desta tarefa termina na exportação dos agrupamentos.

\clearpage
\titulo{4. Agrupamentos no grid}

\begin{center}
\includegraphics[width=.98\linewidth]{../outputs/graficos/grid.png}

\textbf{Figura 4.} Saída do visualizador oficial VictSim3, com sequenciamento desligado.
\end{center}

\textbf{Contornos por cluster.}
\textcolor[RGB]{180,0,180}{1 magenta},
\textcolor[RGB]{0,134,143}{2 ciano},
\textcolor[RGB]{185,85,0}{3 laranja},
\textcolor[RGB]{128,0,255}{4 roxo},
\textcolor[RGB]{0,130,70}{5 verde-menta},
\textcolor[RGB]{195,0,95}{6 rosa}.

\textbf{Círculos por START predito.}
Verde (0), amarelo (1), vermelho (2) e preto (3).
O X branco marca divergência em relação ao gabarito oficial,
consultado apenas na visualização.
A base é a célula ciano preenchida. Células pretas são paredes e os demais tons
distinguem custos de terreno. A origem absoluta está no canto superior esquerdo.

\textbf{Fontes.} Enunciado Tarefa 3, v2 (28/10/2026), relatórios T1 e T2 e
\href{https://github.com/tacla/VictSim3/tree/803a01e54bfa01cde797538d8f2b98d66f12afe2/data_vis_and_stats}{manual e visualizador VictSim3}.
\end{document}
